\documentclass[12pt, a4paper]{article}
\usepackage[utf8]{inputenc}
\usepackage[T1]{fontenc}
\usepackage{amsmath, amssymb, amsthm}
\usepackage{mathrsfs}
\usepackage{hyperref}
\usepackage{geometry}
\usepackage{booktabs}
\usepackage{enumerate}
\geometry{margin=2.5cm}

\theoremstyle{plain}
\newtheorem{theorem}{Teorema}[section]
\newtheorem{lemma}[theorem]{Lema}
\newtheorem{proposition}[theorem]{Proposição}
\newtheorem{corollary}[theorem]{Corolário}
\theoremstyle{definition}
\newtheorem{definition}[theorem]{Definição}
\newtheorem{remark}[theorem]{Observação}
\theoremstyle{remark}
\newtheorem*{note}{Nota}

\title{Motor de Herança Estrutural:\\
	Formalização Algébrica da Fita-Dobra e do\\
	Mecanismo de Promoção de Pares de Goldbach}
\author{Tiago Bandeira\\
	\small\texttt{tiagobandeira.goldbach2025@gmail.com}}
\date{Maio de 2026}

\begin{document}
	
	\maketitle
	
	%--------------------------------------------------------------------
	\begin{abstract}
		%--------------------------------------------------------------------
		Formalizamos o \textbf{Motor de Herança Estrutural}, um mecanismo
		de cobertura dos números pares por somas de dois primos.
		O motor opera sobre dois objectos acoplados: a \textbf{fita-dobra}
		$\mathcal{F}_N^\varepsilon$ (uma matriz $2\times N$ de ímpares em
		percurso zigzag) e a grade $G_{3,C}$ (que fornece eventos restritivos
		via diagonais totalmente primas).
		
		A contribuição geométrica central deste artigo é a observação de que
		a arquitectura em zigzag do acoplamento $\Phi$ já realiza, por
		simetria espacial pura, o pareamento dos extremos da grade com soma
		$2M^+$: não é necessária nenhuma ``dinâmica de movimento'' adicional.
		As janelas $W_i$ são pares antipodais \textbf{fixos} na estrutura da
		grade; o índice $i$ é espacial, não temporal.
		
		A contribuição estrutural central deste artigo é a
		\textbf{separação em duas camadas}:
		\begin{itemize}
			\item \textbf{Camada algébrica} (incondicional): a bijeção de
			acoplamento $\Phi$, definida pela condição $3C=2N$, garante
			por álgebra pura que os extremos opostos de cada janela
			$W_i$ somam directamente $2M^+=4N$, e que a enumeração
			espacial das $N-1$ janelas cobre todos os pares antipodais
			da grade.  O elemento central de toda janela $W_i$ é
			$m^+=2N$, metade de $2M^+$.  Estes resultados independem
			inteiramente de primalidade.
			\item \textbf{Camada aritmética} (condicional): sob a
			\textbf{Hipótese Restritiva}~---~na forma fraca~(HR$^-$),
			que postula a existência de alguma janela $W_i$ em que ambos
			os extremos são primos, ou na forma forte~(HR$^+$), que
			exige uma tripla totalmente prima~---~o motor produz uma
			cadeia certificada de pares de Goldbach.
		\end{itemize}
		
		O \textbf{Teorema do Motor} estabelece, sob HR$^-$, que a cadeia
		$2M_0\to2M_0+2\to\cdots$ cobre todos os pares $\geq2M_0$ como somas
		de dois primos.  HR$^-$ e HR$^+$ são isoladas como as únicas
		hipóteses não provadas e serão o objecto do artigo seguinte, que as
		atacará via o peso oracle-free $\Omega_N^*$~\cite{peso} e a
		desigualdade estrutural~\cite{problemaB}.
	\end{abstract}
	
	\tableofcontents
	
	%====================================================================
	\section{Introdução}
	%====================================================================
	
	A Conjectura de Goldbach afirma que todo número par $>2$ é soma de
	dois primos.  Este trabalho integra-se na
	série~\cite{grade,invariante,peso,problemaB} e formaliza o
	\textbf{Motor de Herança Estrutural} introduzido em~\cite{motor},
	fornecendo as demonstrações algébricas que ali ficaram em aberto.
	
	A ideia central é substituir a pergunta estática
	\begin{center}
		\emph{``$2M$ é soma de dois primos?''}
	\end{center}
	pela pergunta dinâmica
	\begin{center}
		\emph{``dado que $2M$ é soma de dois primos (caso base),\\
			a estrutura garante que $2M+2$ também o é?''}
	\end{center}
	
	O motor opera sobre dois objectos acoplados:
	\begin{enumerate}
		\item A \textbf{fita-dobra} $\mathcal{F}_N^\varepsilon$, que
		organiza os ímpares em percurso zigzag e garante, por
		simetria posição-cega, a existência e a promoção de pares
		simétricos com qualquer soma prescrita.
		\item A \textbf{grade} $G_{3,C}$, que fornece \emph{eventos
			restritivos}~---~configurações totalmente primas numa janela
		$W_i$~---~que certificam a primalidade do par prometido.
	\end{enumerate}
	
	\medskip
	\noindent\textbf{Organização do artigo.}
	As Secções~\ref{sec:fita}--\ref{sec:sigma} constituem a
	\emph{camada algébrica}: todos os resultados são incondicionais e
	independem de primalidade.
	A Secção~\ref{sec:motor} introduz a \emph{camada aritmética}: as
	hipóteses restritivas HR$^-$/HR$^+$, o Motor de Herança e o Teorema
	principal.
	A Secção~\ref{sec:exemplos} ilustra ambas as camadas com exemplos
	numéricos explícitos.
	A Secção~\ref{sec:conclusao} sintetiza a separação e enuncia as
	lacunas em aberto.
	
	%====================================================================
	\section{A Fita-Dobra $\mathcal{F}_N^\varepsilon$}
	\label{sec:fita}
	%====================================================================
	
	\subsection{Definição}
	
	\begin{definition}[Fita-dobra]
		\label{def:fita}
		Fixados $N\geq1$ e $\varepsilon\in\{0,1\}$, a
		\textbf{fita-dobra} $\mathcal{F}_N^\varepsilon$ é a matriz
		$2\times N$ cujas células são:
		\[
		a_k = 2k-1,
		\qquad
		b_k = 4N - 2k + 1 - 2\varepsilon,
		\qquad k=1,\dots,N.
		\]
		A \textbf{linha~1} ($a_k$) percorre os ímpares da esquerda para a
		direita; a \textbf{linha~2} ($b_k$) percorre-os da direita para a
		esquerda (zigzag).  O parâmetro $\varepsilon=0$ é o
		\textbf{modo normal}; $\varepsilon=1$ é o \textbf{modo de
			repetição} (a última coluna da linha~2 repete o valor $b_N-2$,
		``tirando a vaga'' de um ímpar e reduzindo a soma em~$2$).
	\end{definition}
	
	\begin{remark}[Cardinalidade]
		Em ambos os modos, a fita contém exactamente $2N$ células.
		A repetição não altera a cardinalidade~---~substitui um valor por
		outro, mantendo o número de ``quadradinhos'' constante.
	\end{remark}
	
	\subsection{Soma constante (resultado incondicional)}
	
	\begin{proposition}[Soma constante]
		\label{prop:soma}
		Para todo $k\in\{1,\dots,N\}$ e todo $\varepsilon\in\{0,1\}$:
		\[
		a_k + b_k = 4N - 2\varepsilon =: 2M(\varepsilon).
		\]
		Toda coluna da fita soma o mesmo valor $2M(\varepsilon)$,
		independentemente da posição~$k$ e independentemente de
		primalidade.
	\end{proposition}
	
	\begin{proof}
		$a_k+b_k = (2k-1)+(4N-2k+1-2\varepsilon) = 4N-2\varepsilon$.
	\end{proof}
	
	\begin{remark}[Cegueira posicional]
		A Proposição~\ref{prop:soma} é puramente algébrica: a fita é
		\textbf{cega à posição} e \textbf{cega à primalidade}.  Não
		importa em que coluna~$k$ um par $(a_k,b_k)$ se encontra, nem se
		os seus elementos são ou não primos~---~a soma é sempre
		$2M(\varepsilon)$.
	\end{remark}
	
	\subsection{Dois pares consecutivos por fita}
	
	\begin{proposition}[Dois pares por fita]
		\label{prop:dois}
		A mesma fita física $\mathcal{F}_N$ representa simultaneamente
		dois números pares consecutivos:
		\[
		\mathcal{F}_N^0 \;\longrightarrow\; 2M^+ = 4N,
		\qquad
		\mathcal{F}_N^1 \;\longrightarrow\; 2M^- = 4N-2,
		\qquad
		2M^+ - 2M^- = 2.
		\]
	\end{proposition}
	
	\begin{proof}
		Directo da Proposição~\ref{prop:soma} com $\varepsilon=0$ e
		$\varepsilon=1$.
	\end{proof}
	
	\begin{remark}[Razão 2:1]
		Os ímpares estão distribuídos numa densidade que é o dobro da dos
		pares.  Por isso, cada coluna nova adicionada à fita cobre dois
		pares consecutivos, e a alternância de modo
		($\varepsilon=0\leftrightarrow1$) é o mecanismo que acessa cada
		um deles dentro da mesma estrutura física.
	\end{remark}
	
	%====================================================================
	\section{Acoplamento $\Phi$ entre $G_{3,C}$ e $\mathcal{F}_N$}
	\label{sec:phi}
	%====================================================================
	
	\subsection{Condição de acoplamento}
	
	\begin{proposition}[Igualdade de cardinalidade]
		\label{prop:card}
		$G_{3,C}$ e $\mathcal{F}_N$ contêm o mesmo conjunto de ímpares
		se e somente se
		\[
		3C = 2N,
		\]
		ou seja, $N=\frac{3C}{2}$, o que exige $2\mid C$.
	\end{proposition}
	
	\begin{proof}
		$G_{3,C}$ tem $3C$ células e $\mathcal{F}_N$ tem $2N$ células.
		Para que ambas contenham exactamente os mesmos $3C=2N$ ímpares
		consecutivos é necessário e suficiente $3C=2N$.
	\end{proof}
	
	\subsection{Percurso zigzag}
	
	\begin{definition}[Percurso zigzag de $G_{3,C}$]
		\label{def:zigzag-grade}
		O \textbf{percurso zigzag} de $G_{3,C}$ é a sequência
		$\sigma_G:\{1,\ldots,3C\}\to\mathbb{N}_{\text{ímpar}}$ definida
		por:
		\[
		\sigma_G(k) = \begin{cases}
			2k-1
			& k=1,\ldots,C
			\quad\text{(linha~1, }\rightarrow\text{)}\\[4pt]
			2(2C-k)+1
			& k=C+1,\ldots,2C
			\quad\text{(linha~2, }\leftarrow\text{)}\\[4pt]
			2(k-2C)+4C-1
			& k=2C+1,\ldots,3C
			\quad\text{(linha~3, }\rightarrow\text{)}
		\end{cases}
		\]
	\end{definition}
	
	\begin{definition}[Percurso zigzag de $\mathcal{F}_N$]
		\label{def:zigzag-fita}
		O \textbf{percurso zigzag} de $\mathcal{F}_N$ é a sequência
		$\sigma_{\mathcal{F}}:\{1,\ldots,2N\}\to\mathbb{N}_{\text{ímpar}}$
		definida por:
		\[
		\sigma_{\mathcal{F}}(k) = \begin{cases}
			2k-1
			& k=1,\ldots,N
			\quad\text{(linha~1, }\rightarrow\text{)}\\[4pt]
			2(2N-k)+1
			& k=N+1,\ldots,2N
			\quad\text{(linha~2, }\leftarrow\text{)}
		\end{cases}
		\]
	\end{definition}
	
	\subsection{Bijeção de acoplamento (resultado incondicional)}
	
	\begin{definition}[Acoplamento $\Phi$]
		\label{def:phi}
		Dado $C\geq2$ com $2\mid C$, seja $N=\frac{3C}{2}$.
		O \textbf{acoplamento}
		\[
		\Phi: \{1,\ldots,N\} \longrightarrow G_{3,C}\times G_{3,C}
		\]
		associa a cada coluna $k$ da fita $\mathcal{F}_N$ o par de células
		correspondente em $G_{3,C}$, segundo a regra:
		\[
		\Phi(k) = \begin{cases}
			\bigl(f(1,\,k),\; f(3,\,C+1-k)\bigr)
			& k = 1,\ldots,C\\[6pt]
			\bigl(f(2,\,2C+1-k),\; f(2,\,k-C)\bigr)
			& k = C+1,\ldots,\tfrac{3C}{2}
		\end{cases}
		\]
		Para $k\leq C$, a fita emparelha a linha~1 da grade
		(percorrida $\rightarrow$) com a linha~3 (percorrida
		$\leftarrow$).  Para $k>C$, emparelha as extremidades da linha~2
		convergindo para o centro.
	\end{definition}
	
	\begin{proposition}[Acoplamento bem definido e soma constante]
		\label{prop:phi}
		Sob a condição $3C=2N$, $\Phi$ é uma bijeção bem definida.
		Além disso, para todo $k\in\{1,\ldots,N\}$, independentemente
		de primalidade:
		\[
		\Phi(k)_1 + \Phi(k)_2 = 4N = 2M^+.
		\]
	\end{proposition}
	
	\begin{proof}
		\textbf{Bijeção.}  Os pares $\Phi(k)$ para $k=1,\ldots,N$ cobrem
		exactamente os $2N=3C$ ímpares de $G_{3,C}$ sem repetição: para
		$k\leq C$ cobrem as linhas~1 e~3 ($2C$ elementos), e para $k>C$
		cobrem a linha~2 ($C$ elementos, emparelhados das extremidades
		para o centro).  Total: $2C+C=3C=2N$. \checkmark
		
		\textbf{Soma constante para $k\leq C$.}
		\[
		f(1,k)+f(3,C+1-k) = (2k-1)+(4C+2(C+1-k)-1)
		= (2k-1)+(6C-2k+1) = 6C = 4N.
		\]
		
		\textbf{Soma constante para $k=C+j$, $j=1,\ldots,C/2$.}
		Usando $f(2,c)=2(C+c)-1$:
		\[
		f(2,C+1-j)+f(2,j)
		= (4C-2j+1)+(2C+2j-1) = 6C = 4N. \quad\checkmark
		\]
	\end{proof}
	
	\subsection{Par deslocado garantido pelo acoplamento
		(resultado incondicional)}
	
	\begin{proposition}[Par deslocado garantido]
		\label{prop:par-deslocado}
		Seja $W_{i^*}$ a janela central de $G_{3,C}$ com
		$i^*=\lfloor C/2\rfloor$, e seja $(e_1,e_3)$ os extremos de
		qualquer configuração canónica
		$\mathcal{E}\in\{D_1,D_2,C_v\}$ em $W_{i^*}$.
		
		\medskip
		\textbf{Parte~1 --- Deslocamento (incondicional).}
		Via acoplamento $\Phi$, o par $(e_1,e_3)$ satisfaz sempre~---~
		independentemente de primalidade~---:
		\[
		e_1 + e_3 = 2M(\varepsilon) - 2.
		\]
		O par está deslocado de $1$ posição em relação ao par simétrico
		corrente da fita~---~``diagonalizado por $1$ quadradinho''.
		
		\medskip
		\textbf{Parte~2 --- Alinhamento (incondicional).}
		No próximo passo da fita ($\varepsilon:0\to1$ ou $N\to N+1$),
		a soma avança de $2$, independentemente de primalidade:
		\[
		e_1 + e_3 = 2M(\varepsilon)-2
		\;\longrightarrow\;
		e_1 + e_3 = 2M(\varepsilon).
		\]
		O par deslocado torna-se par simétrico \textbf{alinhado} com a
		fita corrente.  Esta é a propriedade geométrica que conecta o
		crescimento da fita ao motor via o passo $2M\to2M+2$.
	\end{proposition}
	
	\begin{proof}
		\textbf{Parte~1.}  Pela Proposição~\ref{prop:phi}, todo par
		$\Phi(k)$ satisfaz $\Phi(k)_1+\Phi(k)_2=4N=2M^+$.  Os extremos
		$(e_1,e_3)$ de $W_{i^*}$ correspondem, via $\Phi$, a um par na
		coluna $k^*$ adjacente ao centro da fita, donde a soma é
		$2M^+-2=2M^-$. \qed
		
		\textbf{Parte~2.}  Por Definição~\ref{def:fita},
		$b_k=4N-2k+1-2\varepsilon$.  O deslocamento mínimo
		($\varepsilon\to0$ ou $N\to N+1$) incrementa $b_k\to b_k+2$ para
		todo~$k$ simultaneamente, logo $e_1+e_3\to e_1+e_3+2$. \qed
	\end{proof}
	
	\begin{remark}[A Parte~3 pertence ao motor, não à álgebra]
		A Proposição~\ref{prop:par-deslocado} garante a existência e o
		alinhamento do par $(e_1,e_3)$ para \emph{qualquer} configuração
		$\mathcal{E}$, prima ou não.  A questão de \emph{quando} esse par
		é de primos~---~e de como isso certifica Goldbach~---~é tratada
		separadamente na Secção~\ref{sec:motor}, sob hipótese aritmética
		explícita.
	\end{remark}
	
	\begin{remark}[Por que fixar o centro]
		A escolha de $i^*=\lfloor C/2\rfloor$ não é arbitrária.  Os
		extremos de $W_{i^*}$ são exactamente os elementos que, via
		$\Phi$, caem na coluna adjacente ao centro da fita~---~garantindo
		o deslocamento de exactamente $2$.  Para janelas $W_i$ com
		$i\neq i^*$, o deslocamento seria diferente e a sincronização com
		o passo da fita quebraria.  O argumento dual~---~por que o
		\emph{passo} de $\sigma$ também deve ser $\delta=1$~---~é
		desenvolvido na Proposição~\ref{prop:pares-diagonais} Parte~3.
	\end{remark}
	
	%====================================================================
	\section{Distribuição Estrutural das Janelas $W_i$}
	\label{sec:sigma}
	%====================================================================
	
	Esta secção é inteiramente incondicional: todos os resultados são
	álgebra pura, sem qualquer hipótese de primalidade.
	
	\subsection{Definição}
	
	\begin{remark}[Natureza espacial do índice $i$]
		\label{rem:espacial}
		A arquitectura em zigzag do acoplamento $\Phi$ realiza, por
		simetria geométrica pura, o pareamento dos extremos opostos da
		grade com soma $2M^+$.  Não há ``deslocamento de peças'': cada
		janela $W_i$ é um par antipodal \textbf{fixo} na estrutura estática
		da grade.  O índice $i \in \{1,\ldots,N-1\}$ é
		\textbf{espacial}~---~identifica qual par antipodal estamos a
		observar~---~e não uma iteração temporal.
		
		A permutação $\sigma$ introduzida abaixo é uma ferramenta de
		enumeração: ela percorre sistematicamente todas as janelas $W_i$
		sem alterar a grade.  O ``passo'' de $\sigma$ é um passo de
		indexação espacial, não de evolução dinâmica.
	\end{remark}
	
	\begin{definition}[Enumeração das janelas $\sigma$]
		\label{def:sigma}
		Seja
		\[
		S = \bigl\{\sigma_G(2),\,\sigma_G(3),\,\ldots,\,
		\sigma_G(3C-1)\bigr\}
		\]
		o conjunto dos $3C-2$ elementos \emph{interiores} do percurso
		zigzag de $G_{3,C}$.  (Excluem-se os dois extremos
		$\sigma_G(1)=1$ e $\sigma_G(3C)=6C-1$.)
		
		A \textbf{enumeração das janelas} $\sigma:S\to S$ é:
		\[
		\sigma\bigl(\sigma_G(k)\bigr) = \sigma_G(k-1),
		\qquad k = 2,\,3,\,\ldots,\,3C-1.
		\]
		Cada passo de $\sigma$ avança o índice espacial de uma posição:
		passa da janela $W_i$ para a janela $W_{i+1}$.  A grade
		permanece inalterada; apenas a nossa ``janela de observação''
		desloca-se.
	\end{definition}
	
	\begin{remark}[Por que excluir os extremos]
		\label{rem:extremos}
		Os extremos $1=\sigma_G(1)$ e $6C-1=\sigma_G(3C)$ são os limites
		do intervalo de ímpares sincronizado com a fita corrente.
		Incluí-los na permutação forçaria $W_{i^*}$ a conter valores fora
		desse intervalo, quebrando a propriedade de deslocamento
		(Proposição~\ref{prop:par-deslocado}).  A exclusão dos extremos
		fecha o ciclo correctamente: após $N-1$ passos de $\sigma$, a
		configuração inicial de $W_{i^*}$ é restituída.
		
		Note-se que esta é a correcção em relação à rotação cíclica $\rho$
		anteriormente proposta: $\rho$ incluía os extremos, quebrando o
		motor; $\sigma$ é restrita a $S$ e é consistente com a fita.
	\end{remark}
	
	\subsection{Cobertura completa dos pares diagonais
		(resultado incondicional)}
	
	\begin{proposition}[Cobertura completa das janelas]
		\label{prop:orbita}
		Seja $i^*=\lfloor C/2\rfloor$ e $k^*$ a coluna correspondente
		na fita.  Para $i=1,\ldots,N-1$, independentemente de primalidade:
		\begin{enumerate}[(i)]
			\item \textbf{Ordenação.}  A enumeração $\sigma$ preserva a
			ordem relativa dos elementos dentro de cada linha da grade.
			\item \textbf{Cobertura.}  A aplicação
			\[
			i \;\mapsto\;
			\text{par }(e_1(i),e_3(i))\text{ de }W_i
			\]
			é uma bijecção de $\{1,\ldots,N-1\}$ sobre o conjunto
			dos $N-1$ pares antipodais interiores da grade $G_{3,C}$.
			\item \textbf{Soma constante $2M^+$ (corrigida).}  Para
			todo $i$, os extremos opostos da janela $W_i$ satisfazem:
			\[
			e_1(i) + e_3(i) = 2M^+ = 4N.
			\]
			Cada janela $W_i$ já aponta directamente para o alvo
			futuro $2M^+$.  Não é necessário nenhum incremento adicional.
			\item \textbf{Invariante do centro.}  O elemento central de
			toda janela $W_i$ é
			\[
			c = \frac{2M^+}{2} = 2N = m^+.
			\]
			A igualdade $e_1(i)+e_3(i) = 2c$ confirma, por álgebra
			pura, que os dois extremos são equidistantes do centro na
			progressão de ímpares.
		\end{enumerate}
	\end{proposition}
	
	\begin{proof}
		\textbf{(i)}  $\sigma$ é uma translação no índice espacial $k$:
		$\sigma(\sigma_G(k))=\sigma_G(k-1)$.  Os elementos de cada linha
		são dispostos em ordem crescente (linhas~1 e~3) ou decrescente
		(linha~2) segundo~$k$; uma translação preserva essa ordem relativa.
		
		\textbf{(ii)}  Antes de qualquer passo, $W_{i^*}$ contém os
		elementos nas posições $i^*,\,i^*+C,\,i^*+2C$ do percurso zigzag.
		Após $i$ passos, esses são substituídos pelos das posições
		$i^*+i,\,i^*+C+i,\,i^*+2C+i$.  À medida que $i$ percorre
		$\{1,\ldots,N-1\}$, essas triplas percorrem exactamente os $N-1$
		pares antipodais interiores de $G_{3,C}$, sem repetição.
		
		\textbf{(iii)}  O par $(e_1(i),e_3(i))$ corresponde, via $\Phi$,
		à coluna $k^*+i$ da fita.  Pela Proposição~\ref{prop:phi}:
		$\Phi(k)_1+\Phi(k)_2 = 4N = 2M^+$ para todo $k$. Portanto
		$e_1(i)+e_3(i) = 4N = 2M^+$.
		
		\textbf{(iv)}  $c = (e_1(i)+e_3(i))/2 = 4N/2 = 2N = m^+$.
	\end{proof}
	
	\begin{corollary}[Cada janela $W_i$ é um par antipodal de soma $2M^+$]
		\label{cor:sigma-desloc}
		Cada incremento do índice espacial $i$ passa para o par antipodal
		seguinte na grade.  A enumeração completa percorre todos os $N-1$
		pares disponíveis.  Este percurso é completamente determinista,
		espacial e independe de primalidade.  Em particular: a grade é,
		ela própria, um ``oráculo'' espacial do alvo $2M^+$~---~as janelas
		não ``chegam'' a $2M^+$ após um incremento; elas \emph{já somam}
		$2M^+$ por construção geométrica.
	\end{corollary}
	
	\begin{remark}[Estrutura direccional da enumeração e natureza dos eventos]
		\label{rem:sigma-direcoes}
		A enumeração $\sigma$ não é apenas uma bijecção abstracta: ela
		impõe uma \emph{estrutura direccional distinta em cada linha}
		da grade, visível na evolução dos extremos de $W_i$ com $i$.
		
		Fixada a janela de índice $i$, os três elementos $(e_1(i),\,p_2(i),\,e_3(i))$
		evoluem com o índice espacial segundo os sentidos do percurso zigzag:
		\begin{itemize}
			\item \textbf{L1} ($\rightarrow$):
			$e_1(i{+}1)=e_1(i)+2$
			--- a janela vê elementos progressivamente maiores de~L1;
			\item \textbf{L2} ($\leftarrow$):
			$p_2(i{+}1)=p_2(i)\pm2$
			--- o sinal depende de $i^*$ estar à esquerda ou à direita
			do centro de~L2;
			\item \textbf{L3} ($\rightarrow$, mas \emph{em sentido contrário}
			ao efeito sobre $W_i$):
			$e_3(i{+}1)=e_3(i)-2$
			--- porque $\Phi$ emparelha~L1 com~L3 em sentidos opostos,
			o elemento de $W_i$ em~L3 decresce enquanto o de~L1 cresce.
		\end{itemize}
		A consequência imediata é:
		\[
		e_1(i{+}1)+e_3(i{+}1)
		=\bigl(e_1(i)+2\bigr)+\bigl(e_3(i)-2\bigr)
		=e_1(i)+e_3(i) = 2M^+,
		\]
		confirmando directamente a constância da soma $2M^+$ da
		Proposição~\ref{prop:orbita}(iii) e revelando a sua
		\emph{razão geométrica}: os dois extremos movem-se em sentidos
		opostos e simétricos, de modo que os incrementos se cancelam.
		
		\medskip
		\noindent\textbf{Nota sobre a natureza dos eventos.}
		Os índices $i=1,\ldots,N-1$ da enumeração das janelas são
		\emph{posições espaciais}~---~todas as configurações possíveis de
		pares antipodais dentro da grade fixa $G_{3,C}$.  Eles não são passos
		temporais do motor: o motor avança com o crescimento da fita
		($N\to N{+}1$), e a grade é actualizada por propagação a cada
		novo passo.  A enumeração $\sigma$ percorre o
		\textbf{espaço de busca} da hipótese restritiva: quais os pares
		$(e_1,e_3)$ disponíveis para o evento~HR$^-$ numa grade fixa.
	\end{remark}
	
	\begin{proposition}[Pares antipodais da grade e unicidade do passo mínimo]
		\label{prop:pares-diagonais}
		Seja $G_{3,C}$ a grade com acoplamento $\Phi$ (condição $3C=2N$).
		Definem-se os \textbf{pares antipodais interiores} de $G_{3,C}$
		como os pares $(e_1(i), e_3(i))$ para $i=1,\ldots,N-1$, onde
		$e_1(i)$ está em L1 e $e_3(i)$ está em L3 em posições
		simétricas segundo $\Phi$.
		
		\medskip
		\noindent\textbf{Parte~1 --- Soma constante $2M^+$ dos pares antipodais
			(incondicional).}
		\[
		e_1(i)+e_3(i) = 4N = 2M^+.
		\]
		Cada par antipodal soma $2M^+$, independentemente de primalidade.
		O centro invariante é $c = m^+ = 2N$.
		
		\medskip
		\noindent\textbf{Parte~2 --- Cobertura bijectiva (incondicional).}
		A aplicação $i\mapsto(e_1(i),e_3(i))$ da
		Proposição~\ref{prop:orbita}(ii) é uma bijecção entre
		$\{1,\ldots,N-1\}$ e o conjunto dos $N-1$ pares antipodais
		interiores de $G_{3,C}$.  Em particular, cada par antipodal
		da grade é \emph{apontado} por $W_i$ em exactamente uma posição
		$i$ da enumeração.
		
		\medskip
		\noindent\textbf{Parte~3 --- Unicidade do passo mínimo $\delta=1$
			(incondicional).}
		Se $\sigma$ fosse substituída por $\sigma^\delta$ com $\delta>1$,
		ocorreriam duas quebras simultâneas:
		\begin{enumerate}[(i)]
			\item \emph{Cobertura parcial}: a enumeração percorreria apenas
			$\lceil(N-1)/\delta\rceil$ dos $N-1$ pares antipodais,
			deixando $N-1-\lceil(N-1)/\delta\rceil>0$ pares sem
			cobertura~---~a bijecção da Parte~2 deixa de ser sobrejectiva.
			\item \emph{Dessincronização com o motor}: cada passo
			corresponderia a um salto de $\delta$ posições em vez
			de~$1$, quebrando a cadeia
			$2M\to2M{+}2$ do Teorema~\ref{thm:motor}.
		\end{enumerate}
		O passo $\delta=1$ é assim a \emph{única} escolha que mantém
		simultaneamente a cobertura bijectiva da Parte~2 e a
		sincronização $N\to N{+}1$ do motor.
	\end{proposition}
	
	\begin{proof}
		\textbf{Parte~1.}  Pela Proposição~\ref{prop:phi},
		$\Phi(k)_1+\Phi(k)_2=4N$ para todo~$k$.\checkmark
		
		\textbf{Parte~2.}  Os $N-1$ pares antipodais correspondem
		bijectivamente às colunas $k=1,\ldots,N-1$ da fita (a
		coluna~$N$ fornece o par central simétrico, não antipodal).  Pela
		Proposição~\ref{prop:orbita}(ii), a enumeração de $\sigma$ cobre
		exactamente essas $N-1$ colunas sem repetição, logo a aplicação
		é uma bijecção.\checkmark
		
		\textbf{Parte~3(i).}  $\sigma^\delta$ tem enumeração de comprimento
		$\lceil(N-1)/\delta\rceil$ sobre as janelas, estritamente menor
		que $N-1$ para $\delta>1$.
		
		\textbf{Parte~3(ii).}  Por
		Observação~\ref{rem:sigma-direcoes}, cada passo de $\sigma$
		avança $e_1$ de~$2$ e recua $e_3$ de~$2$.  Após $\delta$ passos,
		o par aponta a coluna $k^*+\delta$ da fita, saindo da
		sincronização $N\to N{+}1$.\qed
	\end{proof}
	
	\begin{remark}[Fixação canónica da janela e do passo]
		\label{rem:canonico}
		A Observação~\ref{rem:extremos} (por que fixar $i^*$) e a
		Proposição~\ref{prop:pares-diagonais} Parte~3 são complementares:
		a primeira mostra que \emph{alterar a janela} para $i\neq i^*$
		quebra o deslocamento exacto; a segunda mostra que \emph{alterar
			o passo} para $\delta>1$ quebra simultaneamente a cobertura e a
		sincronização.  Juntas, fixam canonicamente tanto a janela
		$W_{i^*}$ como o passo unitário de $\sigma$: são as únicas
		escolhas consistentes com a estrutura do motor.
	\end{remark}
	
	\subsection{Verificação numérica: $G_{3,10}$, $i^*=5$, $N=15$}
	\label{subsec:verif-sigma}
	
	Para $C=10$, $N=15$, $2M^+=60$, $2M^-=58$.
	A janela de índice $i^*=5$ está nas posições $5,\,15,\,25$ do
	percurso zigzag: $\sigma_G(5)=9$, $\sigma_G(15)=31$, $\sigma_G(25)=51$.
	
	\textbf{Nota (correcção):} os extremos $(e_1, e_3)$ de cada janela
	$W_i$ somam $2M^+ = 60$, não $2M^-=58$.  A soma $2M^-$ é o alvo
	da \emph{fita corrente}; a grade aponta geometricamente para o
	alvo \emph{seguinte} $2M^+$.
	
	A tabela abaixo mostra os primeiros índices espaciais~$i$ com os
	pares antipodais $(e_1(i),\,e_3(i))$ de $W_i$.
	A coluna ``par primo?'' é \emph{informação aritmética}~---~não
	interfere na álgebra das colunas anteriores.
	
	\[
	\begin{array}{c|cc|c|c}
		i & e_1(i) & e_3(i) & e_1+e_3 & \text{par primo?}\\
		\hline
		1 & 11 & 49 & 60 & \text{não (49 composto)}\\
		2 & 13 & 47 & 60 & \text{sim}\;\checkmark\\
		3 & 15 & 45 & 60 & \text{não}\\
		4 & 17 & 43 & 60 & \text{sim}\;\checkmark\\
		5 & 19 & 41 & 60 & \text{sim}\;\checkmark\\
		\vdots & \vdots & \vdots & \vdots & \vdots
	\end{array}
	\]
	
	Os extremos $1$ e $6C-1=59$ estão excluídos de $S$ e não entram nas
	janelas~---~confirmando que a enumeração fecha correctamente após
	$N-1=14$ índices.  A soma $e_1(i)+e_3(i)=60=2M^+$ é
	constante para todo~$i$ (Proposição~\ref{prop:orbita}(iii)).
	
	%====================================================================
	\section{O Mecanismo de Promoção}
	\label{sec:promocao}
	%====================================================================
	
	\begin{proposition}[Promoção por deslocamento]
		\label{prop:promocao}
		Seja $\mathcal{F}_N^\varepsilon$ a fita corrente com
		$2M(\varepsilon)=4N-2\varepsilon$, e $(e_1,e_3)$ qualquer par
		simétrico de $\mathcal{F}_N^\varepsilon$ com soma
		$e_1+e_3=2M(\varepsilon)$.  Um \textbf{deslocamento mínimo} é
		qualquer uma das operações:
		\begin{itemize}
			\item \textbf{Alternância de modo}:
			$\varepsilon\to1-\varepsilon$ com $N$ fixo;
			\item \textbf{Nova coluna}: $N\to N+1$ com $\varepsilon$ fixo.
		\end{itemize}
		
		\medskip
		\textbf{Parte~1 --- Fita (incondicional).}
		O deslocamento mínimo incrementa $b_k\to b_k+2$ para todo~$k$
		simultaneamente~---~independentemente de primalidade.  Logo, para
		qualquer par $(e_1,e_3)$, primo ou não:
		\[
		e_1 + e_3^{\text{novo}}
		= e_1 + (e_3+2)
		= 2M(\varepsilon)+2.
		\]
		
		\medskip
		\textbf{Parte~2 --- Grade (incondicional).}
		O deslocamento mínimo propaga-se para $G_{3,C}$ como uma
		\textbf{operação de pilha global}: fita e grade partilham os
		mesmos $3C=2N$ ímpares na mesma ordem, pelo que todo elemento de
		$G_{3,C}$ é reindexado simultaneamente.  Após o deslocamento,
		$f'(r,c) = f(r,c) + 2\delta(r)$, onde:
		\[
		\text{Alternância de modo:}\quad
		\delta(1)=0,\;\delta(2)=2,\;\delta(3)=0;
		\]
		\[
		\text{Nova coluna:}\quad
		\delta(1)=0,\;\delta(2)=0,\;\delta(3)=2.
		\]
		Em consequência, toda janela $W_i$ é reorganizada com novos
		valores~---~independentemente de primalidade.
	\end{proposition}
	
	\begin{proof}
		\textbf{Parte~1.}
		Por Definição~\ref{def:fita}, $b_k=4N-2k+1-2\varepsilon$.
		\begin{itemize}
			\item Alternância $\varepsilon\to0$ (se $\varepsilon=1$):
			$b_k^{\text{novo}}=4N-2k+1=b_k+2$. \checkmark
			\item Nova coluna $N\to N+1$ com $\varepsilon\to1$:
			$b_k^{\text{novo}}=4(N+1)-2k+1-2=b_k+4-2=b_k+2$. \checkmark
		\end{itemize}
		
		\textbf{Parte~2.}
		Os $3C=2N$ ímpares de $G_{3,C}$ e de $\mathcal{F}_N$ são os
		mesmos, dispostos em percursos zigzag compatíveis (Secção~\ref{sec:phi}).
		O deslocamento introduz um novo ímpar no extremo superior
		(ou repete um), empurrando todos os outros uma posição na ordem.
		Como a grade preenche as células linha a linha em ordem crescente,
		o empurrão global redistribui os valores por todas as linhas e
		colunas, reorganizando toda janela $W_i$. \qed
	\end{proof}
	
	\begin{remark}[O mecanismo é uma ferramenta algébrica]
		A Proposição~\ref{prop:promocao} mostra que a fita e a grade
		têm um mecanismo de deslocamento completamente determinista.
		Esse mecanismo funciona para \emph{qualquer} par $(e_1,e_3)$,
		seja ele de primos ou não.  O motor de herança usa esta ferramenta
		sob uma hipótese aritmética adicional (Secção~\ref{sec:motor}):
		a de que existe algum passo $t$ em que o par é efectivamente de
		primos.
	\end{remark}
	
	%====================================================================
	\section{Par Base e Par Prometido}
	\label{sec:pares}
	%====================================================================
	
	Esta secção introduz as definições que ligam a camada algébrica ao
	motor de herança.  O par base e o par prometido são definidos em
	termos de primalidade, mas a sua existência como pares simétricos
	deslocados é garantida incondicionalmente pelas secções anteriores.
	
	\begin{definition}[Par base]
		\label{def:base}
		Um par $(a_s,b_s)$ de $\mathcal{F}_N^1$ é um \textbf{par base}
		para $2M^-$ se $a_s$ e $b_s$ são ambos primos.  A escolha da
		coluna $s$ é livre~---~independente de qualquer evento restritivo.
		A existência de um par base é a afirmação de que $2M^-$ é soma
		de dois primos (caso base de Goldbach); o par base não é
		construído pelo motor, mas é o ponto de partida a partir do qual
		o motor opera.
	\end{definition}
	
	\begin{definition}[Par prometido]
		\label{def:prometido}
		O \textbf{par prometido} para $2M^+$ é um par $(p_1,p_3)$ de
		$\mathcal{F}_N^0$ tal que:
		\begin{enumerate}[(i)]
			\item $p_1+p_3=2M^+=4N$
			(garantido pela álgebra do acoplamento $\Phi$~---~incondicional);
			\item $(p_1,p_3)$ são os extremos de $W_{i^*}$ para algum
			passo $t$ da permutação deslizante $\sigma$
			(a associação com $W_{i^*}$ é incondicional~---~
			Proposição~\ref{prop:orbita});
			\item $p_1$ e $p_3$ são ambos primos
			(esta é a condição aritmética, postulada pela Hipótese
			Restritiva~HR$^-$~---~condicional);
			\item $(p_1,p_3)$ estão em uma das diagonais secundárias de $\mathcal{F}_N^0$.
		\end{enumerate}
		Quando existe, o par prometido é o certificado de Goldbach para
		$2M^+$: após o deslocamento mínimo da fita (passo $+2$), ele
		torna-se o par base para $2M^++2$.
	\end{definition}
	
	\begin{remark}[Separação das condições]
		As condições~(i) e~(ii) da Definição~\ref{def:prometido} são
		incondicionais: qualquer par de extremos de $W_{i^*}$ sob
		$\sigma^t$ satisfaz automaticamente $p_1+p_3=2M^+$ e está
		associado a $W_{i^*}$ via a órbita de $\sigma$.  Só a
		condição~(iii)~---~a primalidade simultânea de $p_1$ e $p_3$~---~é
		hipotética.  Esta separação é a essência da arquitectura do motor.
	\end{remark}
	
	\begin{proposition}[Independência entre par base e par prometido]
		\label{prop:indep}
		O par base $(a_s,b_s)\in\mathcal{F}_N^1$ e o par prometido
		$(p_1,p_3)\in\mathcal{F}_N^0$ são \textbf{independentes}: cada
		um soma o respectivo $2M(\varepsilon)$ por
		Proposição~\ref{prop:soma}, mas são determinados por mecanismos
		distintos~---~o par base por qualquer coluna prima da fita, o par
		prometido pelo evento restritivo de $W_{i^*}$.  Em particular,
		não precisam partilhar elementos.
	\end{proposition}
	
	\begin{proof}
		O par base existe em $\mathcal{F}_N^1$ ($\varepsilon=1$) e o par
		prometido em $\mathcal{F}_N^0$ ($\varepsilon=0$); os dois modos
		são distintos e independentes por construção.
	\end{proof}
	
	%====================================================================
	\section{O Motor de Herança: Hipóteses e Teorema}
	\label{sec:motor}
	%====================================================================
	
	As Secções~\ref{sec:fita}--\ref{sec:pares} estabeleceram a camada
	algébrica: toda a geometria dos pares deslocados, a cobertura pela
	órbita de $\sigma$, e o mecanismo de promoção são incondicionais.
	Esta secção introduz a camada aritmética: as hipóteses que postulam
	a existência de pares de primos, e o teorema que delas decorre.
	
	\subsection{As duas hipóteses restritivas}
	
	\begin{definition}[Hipóteses Restritivas]
		\label{def:HR}
		Fixados $C$ e $i^*=\lfloor C/2\rfloor$, sejam
		$(e_1(i),\,e_3(i))$ os extremos da janela $W_i$.
		
		\medskip
		\textbf{Hipótese Restritiva Fraca~(HR$^-$).}
		\[
		\exists\, i\in\{1,\ldots,N-1\}:
		\quad e_1(i)\in\mathbb{P}
		\;\;\text{e}\;\;
		e_3(i)\in\mathbb{P}.
		\]
		Em alguma janela $W_i$, \emph{ambos} os extremos são primos.
		O par $(e_1(i),e_3(i))$ satisfaz automaticamente
		$e_1(i)+e_3(i)=2M^+$ por construção geométrica; esta é a
		condição mínima para o motor.
		
		\medskip
		\textbf{Hipótese Restritiva Forte~(HR$^+$).}
		\[
		\exists\, i\in\{1,\ldots,N-1\}:
		\quad \bigl(e_1(i),\,p_2(i),\,e_3(i)\bigr)
		\text{ é totalmente prima.}
		\]
		Existe uma janela $W_i$ em que a \emph{tripla completa}
		(diagonal $D_1$, $D_2$ ou coluna central $C_v$) é inteiramente
		formada por primos.  HR$^+$ implica HR$^-$.
	\end{definition}
	
	\begin{remark}[Relação entre HR$^-$ e HR$^+$]
		HR$^+$ implica HR$^-$ (uma tripla totalmente prima fornece em
		particular dois extremos primos).  O recíproco é falso em geral.
		O motor opera sob HR$^-$; HR$^+$ é a ferramenta para
		demonstrações mais profundas~---~nomeadamente a prova da própria
		HR$^-$ via densidade de eventos restritivos e a recorrência da
		cadeia de Markov associada.
	\end{remark}
	
	\subsection{Teorema do Motor de Herança}
	
	\begin{theorem}[Motor de Herança Estrutural]
		\label{thm:motor}
		\textit{Hipóteses:}
		\begin{itemize}
			\item[\textbf{(HI)}] Existe $2M_0$ com par base
			$(a_s,b_s)\in\mathcal{F}_{N_0}^1$ (caso base de Goldbach)
			e evento restritivo inicial em $G_{3,C_0}$.
			\item[\textbf{(HR$^-$)}] Para toda iteração do motor, existe
			pelo menos uma janela $W_i$ ($i\in\{1,\ldots,N-1\}$) tal que
			$(e_1(i),e_3(i))$ são ambos primos.
			A soma $e_1(i)+e_3(i)=2M^+$ é garantida geometricamente.
		\end{itemize}
		
		\textit{Conclusão.}  A cadeia
		\[
		2M_0 \;\longrightarrow\; 2M_0+2 \;\longrightarrow\; 2M_0+4
		\;\longrightarrow\; \cdots
		\]
		cobre todos os pares $\geq 2M_0$ como somas de dois primos.
	\end{theorem}
	
	\begin{proof}
		Procedemos por indução, usando a camada algébrica das
		Secções~\ref{sec:fita}--\ref{sec:pares} em cada passo.
		
		\textbf{Base.}  Por~HI, $2M_0$ é Goldbach via par base
		$(a_s,b_s)\in\mathcal{F}_{N_0}^1$.  Por~HR$^-$, existe $i_0$
		tal que $(e_1(i_0),e_3(i_0))$ são ambos primos; pela
		Definição~\ref{def:prometido}, este par é o par prometido
		$(p_1,p_3)$ com $p_1+p_3=2M^+=2M_0+2$.
		
		\textbf{Promoção.}  Na iteração seguinte $\mathcal{F}_{N_1}^1$,
		o par $(p_1,p_3)$ é promovido a par base para $2M_0+2$
		pela Proposição~\ref{prop:promocao} Parte~1~---~certificando
		Goldbach para $2M_0+2$.
		
		\textbf{Regeneração.}  Por~HR$^-$ aplicada à nova configuração,
		existe $i_1$ tal que $(e_1(i_1),e_3(i_1))$ são ambos primos;
		este par é o novo par prometido $(q_1,q_3)$ com
		$q_1+q_3=2M_0+4$.
		
		\textbf{Indução.}  Repetindo promoção e regeneração, cada par
		$2M_0+2n$ é certificado como soma de dois primos.
	\end{proof}
	
	\begin{remark}[Estatuto de cada componente]
		\label{rem:estatuto}
		\begin{center}
			\begin{tabular}{lll}
				\toprule
				Componente & Estatuto & Referência \\
				\midrule
				Soma constante da fita & Incondicional &
				Prop.~\ref{prop:soma} \\
				Bijeção $\Phi$ e soma constante & Incondicional &
				Prop.~\ref{prop:phi} \\
				Par deslocado garantido (Partes~1--2) & Incondicional &
				Prop.~\ref{prop:par-deslocado} \\
				Órbita de $\sigma$ (cobertura, soma) & Incondicional &
				Prop.~\ref{prop:orbita} \\
				Promoção por deslocamento (Partes~1--2) & Incondicional &
				Prop.~\ref{prop:promocao} \\
				HI (caso base) & Verificável &
				Green--Tao~\cite{green-tao} garante existência \\
				HR$^-$ (par de primos em $W_{i^*}$) & Condicional &
				Lacuna central (mínima para o motor) \\
				HR$^+$ (tripla totalmente prima) & Condicional &
				Versão forte; implica HR$^-$ \\
				Teorema~\ref{thm:motor} & Condicional a HR$^-$ & --- \\
				\bottomrule
			\end{tabular}
		\end{center}
	\end{remark}
	
	\begin{remark}[Cadeia de Markov cíclica]
		O motor com a permutação $\sigma$ tem a estrutura de uma
		\textbf{cadeia determinista cíclica}: o estado é o par
		$(2M,\,\mathcal{E}(t))$ e a transição $\sigma$ é determinística
		na álgebra.  HR$^-$ é equivalente à \textbf{recorrência} da
		cadeia~---~a cadeia não tem estados absorventes, i.e., nenhum
		par $2M$ é evitado para sempre.  Esta conexão com teoria
		ergódica~\cite{grade} será explorada no artigo seguinte.
	\end{remark}
	
	%====================================================================
	\section{Exemplos Numéricos}
	\label{sec:exemplos}
	%====================================================================
	
	\begin{remark}[Condição de acoplamento nos exemplos]
		O acoplamento $\Phi$ exige $3C=2N$, logo $3\mid N$.  O menor
		exemplo não trivial é $C=8$, $N=12$
		($3\times8=2\times12=24$).  Exemplos com $N$ não divisível por
		$3$ (e.g.\ $N=11$) descrevem a fita isoladamente, sem acoplamento
		a uma grade inteira.
	\end{remark}
	
	\subsection{Fita isolada: $N=11$, $2M^+=44$, $2M^-=42$}
	
	\[
	\begin{array}{c|ccccccccccc}
		k    & 1 & 2 & 3 & 4 & 5 & 6 & 7 & 8 & 9 & 10 & 11\\
		\hline
		a_k  & 1 & 3 & 5 & 7 & 9 &11 &13 &15 &17 &19  &21 \\
		b_k^0& 43&41 &39 &37 &35 &33 &31 &29 &27 &25  &23 \\
		b_k^1& 41&39 &37 &35 &33 &31 &29 &27 &25 &23  &21 \\
	\end{array}
	\]
	
	\begin{itemize}
		\item \textbf{Álgebra (incondicional):} toda coluna soma $44$
		em modo $\varepsilon=0$ e $42$ em modo $\varepsilon=1$,
		independentemente de primalidade.
		\item \textbf{Par base (aritmético):}
		$\mathcal{F}_{11}^0$ tem par primo $(13,31)$ na coluna
		$k=7$ ($13+31=44$~\checkmark);
		$\mathcal{F}_{11}^1$ tem par primo $(13,29)$ na coluna
		$k=7$ ($13+29=42$~\checkmark).
	\end{itemize}
	
	\subsection{Acoplamento completo: $C=8$, $N=12$, $2M^+=48$,
		$2M^-=46$}
	
	$3\times8=2\times12=24$ células. \checkmark
	
	\medskip\noindent\textbf{Grade $G_{3,8}$:}
	\[
	\begin{array}{c|cccccccc}
		& c=1&2&3&4&5&6&7&8\\
		\hline
		\text{Linha~1 }(\rightarrow) & 1&3&5&7&9&11&13&15\\
		\text{Linha~2 }(\leftarrow)  &31&29&27&25&23&21&19&17\\
		\text{Linha~3 }(\rightarrow) &33&35&37&39&41&43&45&47\\
	\end{array}
	\]
	
	\medskip\noindent\textbf{Fita $\mathcal{F}_{12}$:}
	\[
	\begin{array}{c|cccccccccccc}
		k & 1&2&3&4&5&6&7&8&9&10&11&12\\
		\hline
		a_k   &1&3&5&7&9&11&13&15&17&19&21&23\\
		b_k^0 &47&45&43&41&39&37&35&33&31&29&27&25\\
		b_k^1 &45&43&41&39&37&35&33&31&29&27&25&23\\
	\end{array}
	\]
	
	\medskip\noindent\textbf{Acoplamento $\Phi$~---~camada algébrica
		(incondicional):}
	\[
	\begin{array}{c|cc|c|c}
		k & \Phi(k)_1 & \Phi(k)_2 & \text{soma} & \text{células em }G_{3,8}\\
		\hline
		1  & 1  & 47 & 48 & f(1,1)+f(3,8)\\
		2  & 3  & 45 & 48 & f(1,2)+f(3,7)\\
		3  & 5  & 43 & 48 & f(1,3)+f(3,6)\\
		4  & 7  & 41 & 48 & f(1,4)+f(3,5)\\
		5  & 9  & 39 & 48 & f(1,5)+f(3,4)\\
		6  & 11 & 37 & 48 & f(1,6)+f(3,3)\\
		7  & 13 & 35 & 48 & f(1,7)+f(3,2)\\
		8  & 15 & 33 & 48 & f(1,8)+f(3,1)\\
		9  & 17 & 31 & 48 & f(2,8)+f(2,1)\\
		10 & 19 & 29 & 48 & f(2,7)+f(2,2)\\
		11 & 21 & 27 & 48 & f(2,6)+f(2,3)\\
		12 & 23 & 25 & 48 & f(2,5)+f(2,4)\\
	\end{array}
	\]
	Toda coluna soma $48=4N=2M^+$, independentemente de primalidade. \checkmark
	
	\medskip\noindent\textbf{Motor em acção~---~camada aritmética
		(condicional a HR$^-$):}
	\begin{itemize}
		\item $\mathcal{F}_{12}^1$ ($2M^-=46$): par base primo $(5,41)$
		na coluna $k=3$ ($5+41=46$~\checkmark).
		\item \textbf{Janelas da grade $G_{3,8}$ (corrigido):}
		As janelas $W_i$ entregam pares com soma $2M^+=48$ directamente,
		por construção geométrica.  Os primeiros pares antipodais são:
		$W_1=(1,47)$, $W_2=(3,45)$, $W_3=(5,43)$, $W_4=(7,41)$,
		$W_5=(11,37)$, \ldots~---~todos somam $48$~\checkmark.
		O centro invariante de todas as janelas é $c=24=m^+$~\checkmark.
		\item \textbf{Evento HR$^-$:} $W_4=(7,41)$, ambos primos
		($7+41=48=2M^+$~\checkmark).  Também $W_5=(11,37)$: $11+37=48$~\checkmark.
		\item \textbf{Promoção (incondicional):} o par $(7,41)$ torna-se
		par base para $2M^+=48$ pelo mecanismo de
		Proposição~\ref{prop:promocao}~---~Goldbach para
		$48$~\checkmark.
		\item Par prometido para $2M=50$: determinado pelo próximo
		evento restritivo em $G_{3,C'}$ após deslocamento mínimo
		(requer HR$^-$ na nova configuração).
	\end{itemize}
	
	%====================================================================
	\section{Conclusão}
	\label{sec:conclusao}
	%====================================================================
	
	Este artigo constitui a \textbf{fundação algébrica} do Motor de
	Herança Estrutural na série sobre a Conjectura de Goldbach.
	
	\medskip
	\noindent\textbf{Camada algébrica~---~resultados incondicionais.}
	Demonstrámos por álgebra pura, sem qualquer hipótese de primalidade:
	\begin{enumerate}
		\item A \textbf{soma constante} $a_k+b_k=2M(\varepsilon)$ em
		toda coluna da fita~---~propriedade posição-cega e
		primalidade-cega (Proposição~\ref{prop:soma}).
		\item A \textbf{bijeção} $\Phi$ entre $G_{3,C}$ e
		$\mathcal{F}_N$, com soma constante $\Phi(k)_1+\Phi(k)_2
		=4N$ para todo~$k$ (Proposição~\ref{prop:phi}).
		\item O \textbf{par deslocado garantido}: os extremos de
		$W_{i^*}$ somam sempre $2M(\varepsilon)-2$, e o próximo
		deslocamento mínimo os alinha com $2M(\varepsilon)$~---~
		quer sejam ou não primos
		(Proposição~\ref{prop:par-deslocado}).
		\item A \textbf{distribuição estrutural das janelas}: os $N-1$
		pares antipodais são cobertos em bijecção pela enumeração
		espacial $\sigma$, com soma constante $2M^+$ em toda janela
		e elemento central invariante $c=m^+=2N$
		(Proposição~\ref{prop:orbita}).  \textbf{[Corrigido:} a soma
		constante é $2M^+$, não $2M^-$~---~a grade é um ``oráculo''
		espacial do alvo futuro.\textbf{]}
		\item A \textbf{estrutura direccional} de $\sigma$ e a
		\textbf{identidade dos pares diagonais}: L1 e L3
		movem-se em sentidos opostos ($e_1\!\uparrow$,
		$e_3\!\downarrow$), os extremos de $W_{i^*}$ apontam
		exactamente os pares $(a_k,b_{k+1}^{(0)})$ com soma
		$2M^-$, e o passo unitário $\delta=1$ é a única escolha
		que mantém simultaneamente a cobertura bijectiva e a
		sincronização $N\to N{+}1$
		(Observação~\ref{rem:sigma-direcoes};
		Proposição~\ref{prop:pares-diagonais}).
		\item O \textbf{mecanismo de promoção}: qualquer deslocamento
		$+2$ na fita propaga-se para $G_{3,C}$ como operação de
		pilha global, promovendo qualquer par simétrico~---~primo
		ou não~---~ao alinhamento com a fita seguinte
		(Proposição~\ref{prop:promocao}).
	\end{enumerate}
	
	\medskip
	\noindent\textbf{Camada aritmética~---~resultados condicionais.}
	Sob as Hipóteses Restritivas HR$^-$/HR$^+$, que postulam a
	existência de pelo menos um passo $t$ da órbita de $\sigma$ em que
	os extremos de $W_{i^*}$ são ambos primos (HR$^-$) ou a tripla
	completa é totalmente prima (HR$^+$):
	\begin{enumerate}
		\item O \textbf{par prometido} existe
		(Definição~\ref{def:prometido}): os seus extremos, cujo
		alinhamento é garantido algebricamente, são de primos por
		hipótese.
		\item O \textbf{Teorema do Motor}
		(Teorema~\ref{thm:motor}) estabelece que a cadeia
		$2M_0\to2M_0+2\to\cdots$ cobre todos os pares $\geq 2M_0$
		como somas de dois primos.
	\end{enumerate}
	
	\medskip
	\noindent\textbf{Lacuna remanescente.}
	A única hipótese não provada é a primalidade simultânea dos
	extremos $(e_1(t),e_3(t))$ em algum passo~$t$~---~i.e., HR$^-$.
	A existência geométrica do par e a sua cobertura pela órbita de
	$\sigma$ são incondicionais; a primalidade é o único ingrediente
	aritmético.  HR$^-$ (e a versão mais forte HR$^+$) serão o objecto
	do artigo seguinte, que as atacará usando o peso oracle-free
	$\Omega_N^*$~\cite{peso} e a desigualdade estrutural
	$\mu_{\text{comp}}>\mu_{\text{global}}>\mu_{\text{primo}}$~\cite{problemaB},
	com o acoplamento $\Phi$ e a órbita de $\sigma$ disponíveis como
	ferramentas algébricas.
	
	%====================================================================
	\begin{thebibliography}{9}
		%====================================================================
		
		\bibitem{grade}
		T.~Bandeira,
		\textit{Uma Abordagem Unificada para a Conjectura de Goldbach:
			Estrutura Combinatória, Saturação Geométrica e o Elo entre Trios
			e Pares Primos}, preprint, maio de 2026.
		
		\bibitem{invariante}
		T.~Bandeira,
		\textit{Uma propriedade de soma constante na grade $3\times N$
			e as conjecturas de Goldbach}, preprint, maio de 2026.
		
		\bibitem{peso}
		T.~Bandeira,
		\textit{O Invariante Âncora da Grade $G_{3,N}$ e um Estimador
			Geométrico Oracle-Free para a Conjectura de Goldbach},
		preprint, maio de 2026.
		
		\bibitem{problemaB}
		T.~Bandeira,
		\textit{Problema~B do Crivo Geométrico: Demonstração da
			Desigualdade Estrutural}, preprint, maio de 2026.
		
		\bibitem{motor}
		T.~Bandeira,
		\textit{Motor de Herança Estrutural em Grades $3\times N$:
			Um Sistema Dinâmico de Cobertura Condicional para Pares de
			Goldbach}, preprint, maio de 2026.
		
		\bibitem{green-tao}
		B.~Green e T.~Tao,
		\textit{The primes contain arbitrarily long arithmetic
			progressions},
		\textit{Annals of Mathematics}, vol.~167, pp.~481--547, 2008.
		
	\end{thebibliography}
	
\end{document}